\ChapterImageStar[cap:marcoConceptual]{Marco Conceptual}{./images/fondo.png}\label{cap:marcoConceptual}
\mbox{}\\
%------------------------------------------------------------------------------------------------------------------------------------
\noindent
En esta sección se presentan los conceptos fundamentales que permiten comprender el contexto de la propuesta, enfocados en la 
computación en la nube y la seguridad de la información en entornos tecnológicos.

\noindent
\section{Computación en la nube}
\noindent
La computación en la nube se define como un modelo que permite el acceso universal, conveniente y bajo demanda a recursos computacionales 
compartidas, tales como redes, servidores, almacenamiento y servicios. En este sentido, según~\cite{mell2011nist}, este modelo se 
caracteriza por ofrecer servicios de manera flexible, escalable y con un mínimo esfuerzo de gestión o interacción con el proveedor del servicio.

Adicionalmente, este paradigma se fundamenta en cinco caracteristicas principales autoservicio bajo demanda, amplio acceso a la red, agrupación 
de recursos, rápida elasticidad y servicio medio. Asimismo, contempla diferentes modelos de servicio, entre los cuales se encuentran infraestructura 
como servicio (\IaaS), plataforma como servicio (\PaaS) y software como servicio (\SaaS), así como modelos de despliegue que incluyen nubes públicas,
privadas, híbridas y comunitarias.

\section{Tríada \CIA}
\noindent
La tríada \CIA \textit{(Confidentiality, Integrity, Availability)} es un modelo que define los principios fundamentales de la seguridad de 
la información: confidencialidad, integridad y disponibilidad.~\cite{stallings2012computer}

\section{Seguridad en la nube}
\noindent
En este contexto, la seguridad en la nube se entiende como el conjunto de políticas, tecnologías y controles destinados a proteger 
los activos de información en entornos de computación en la nube. La \textit{European Network and Information Security Agency}~\cite{enisa2009cloud}, 
establece que esta debe evaluarse considerando los riesgos asociados a la confidencialidad, integridad y disponibilidad de la información.

\section{Seguridad de la información}
\noindent
Por su parte, la seguridad de la información se orienta a garantizar la protección de los datos frente a accesos no autorizados,
alteraciones indebidas o pérdida de disponibilidad, asegurando así los principios de confidencialidad, integridad y disponibilidad. 
De acuerdo con la norma ISO/IEC 27001:2022, la gestión de la seguridad de la información debe realizarse mediante un enfoque 
sistemático basado en la identificación, análisis y tratamiento de riesgos, a través de la implementación de un Sistema de Gestión 
de Seguridad de la Información (\SGSI), que permita establecer, operar, monitorear, revisar y mejorar continuamente los controles 
de seguridad.~\cite{iso27001}

\section{Activos de información}
\noindent
Finalmente, los activos de información corresponden a todos aquellos recursos que poseen valor para una organización, incluyendo 
datos, sistemas, servicios, infraestructura tecnológica y recursos humanos asociados al procesamiento de la información. De acuerdo a 
la norma ISO/IEC 27001:2022, estos activos deben ser identificados, clasificados y protegidos mediante controles adecuados, en función
de su valor, criticidad y nivel de riesgo, con el fin de garantizar su seguridad a lo largo de su ciclo de vida.

En entornos de computación en la nube, la protección de estos activos adquiere especial relevancia, particularmente cuando involucran
información sensible o datos personales. Según la norma ISO/IEC 27018:2025 establece lineamientos específicos para la Protección de 
Información de Identificación personal (\PII), en servicios de nube publica, fortaleciendo la confidencialidad y privacidad de los
datos.~\cite{iso27018}

